\section{实验环节：面部检测 Lab}

Amini 在课堂上介绍了本讲配套的实验（Lab 2）——面部检测任务。这个实验让学生亲手构建一个 CNN 来检测图像中的人脸，是将本讲理论付诸实践的绝佳机会。

实验的主要步骤包括：
\begin{enumerate}
    \item 构建卷积神经网络架构（包括卷积层、ReLU 激活、池化层和全连接分类头）
    \item 使用面部图像数据集训练模型
    \item 评估模型在新图像上的检测性能
    \item 分析模型学到的特征可视化
\end{enumerate}

\begin{knowledgebox}{面部检测的历史意义}
面部检测是最早获得大规模商业应用的计算机视觉任务之一。从 Viola-Jones 人脸检测器（2001年，基于手工设计的 Haar 特征和 AdaBoost 分类器）到现代的深度学习方法，面部检测的精度和速度都有了质的飞跃。如今，面部检测技术已经深入到手机解锁、社交媒体滤镜、安防监控等日常应用中。
\end{knowledgebox}

\begin{warningbox}{面部识别的伦理考量}
虽然面部检测和识别技术非常强大，但也引发了严重的隐私和伦理问题。包括：种族和性别偏见（模型在某些群体上表现更差）、监控滥用、Deepfake 伪造等。在开发和部署这类技术时，必须认真考虑其社会影响。
\end{warningbox}

\newpage
